\section{Конструкторский раздел}

В данном разделе приводится описание структуры разрабатываемой базы данных. В
соответствии с аналитическим разделом была выбрана гибридная модель хранения данных.
Для реляционной базы данных каждой выделенной сущности ставится в соответствие таблица, а
для графовой базы данных описываются структуры узлов и ребер. Рассматриваются
спроектированные ограничения целостности, а также ролевая модель управления доступом на
уровне приложения.

\subsection{Проектирование базы данных}

На основе формализации данных, проведенной в аналитическом разделе, в разрабатываемой
реляционной базе данных выделены следующие таблицы:
\begin{itemize}
    \item teams --- таблица команд разработчиков;
    \item users --- таблица пользователей приложения;
    \item invites --- таблица приглашений для регистрации;
    \item it\_systems --- таблица систем.
\end{itemize}

В таблицах \ref{tbl:teams} --- \ref{tbl:it_system} определены типы данных и ограничения
для каждого столбца перечисленных реляционных таблиц.

\begin{center}
    \begin{tabularx}{\textwidth}{|p{2.4cm}|X|p{3cm}|X|}
        \caption{\label{tbl:teams}Информация о таблице teams} \\ \hline
        Столбец & Значение & Тип & Ограничение \\ \hline
        \endfirsthead

        \multicolumn{4}{l}{{Продолжение таблицы\ \thetable{}}} \\ \hline
        Столбец & Значение & Тип & Ограничение \\ \hline
        \endhead

        \endfoot
        \endlastfoot

        id & Идентификатор & UUID & Не null, первичный ключ \\ \hline
        name & Название & Строка & Не null, уникальный \\ \hline
        description & Описание & Строка & \\ \hline
    \end{tabularx}
\end{center}

\begin{center}
    \begin{tabularx}{\textwidth}{|p{3cm}|X|p{3cm}|X|}
        \caption{\label{tbl:user_account}Информация о таблице users} \\ \hline
        Столбец & Значение & Тип & Ограничение \\ \hline
        \endfirsthead

        \multicolumn{4}{l}{{Продолжение таблицы\ \thetable{}}} \\ \hline
        Столбец & Значение & Тип & Ограничение \\ \hline
        \endhead

        \endfoot
        \endlastfoot

        id & Идентификатор & UUID & Не null, первичный ключ \\ \hline
        username & Имя пользователя & Строка & Не null, уникальный \\ \hline
        email & Электронная почта & Строка & Не null, уникальный \\ \hline
        password\_hash & Хэш пароля & Строка & Не null \\ \hline
        role & Идентификатор \newline роли & Целое число & Не null \\ \hline
        team\_id & Идентификатор \newline команды & UUID & Внешний ключ \\ \hline
    \end{tabularx}
\end{center}
\begin{center}
    \begin{tabularx}{\textwidth}{|X|X|p{3cm}|X|}
        \caption{\label{tbl:invites}Информация о таблице invites} \\ \hline
        Столбец & Значение & Тип & Ограничение \\ \hline
        \endfirsthead

        \multicolumn{4}{l}{{Продолжение таблицы\ \thetable{}}} \\ \hline
        Столбец & Значение & Тип & Ограничение \\ \hline
        \endhead

        \endfoot
        \endlastfoot

        id & Идентификатор & UUID & Не null, \newline первичный ключ \\ \hline
        status & Статус приглашения & Целое число & Не null \\ \hline
        target\_email & Электронная \newline почта получателя & Строка & Не null, \newline уникальный \\ \hline
        role & Роль нового \newline пользователя & Целое число & Не null \\ \hline
        code & Уникальный код & Строка & Не null, \newline уникальный \\ \hline
        expiration\_date & Дата истечения & Дата и время & Не null \\ \hline
        team\_id & Идентификатор команды & UUID & Не null, внешний ключ \\ \hline
        activated\_by\_uid & Идентификатор нового \newline пользователя & UUID & Не null, внешний ключ \\ \hline
    \end{tabularx}
\end{center}

\begin{center}
    \begin{tabularx}{\textwidth}{|p{3cm}|X|p{3cm}|X|}
        \caption{\label{tbl:it_system}Информация о таблице it\_systems} \\ \hline
        Столбец & Значение & Тип & Ограничение \\ \hline
        \endfirsthead

        \multicolumn{4}{l}{{Продолжение таблицы\ \thetable{}}} \\ \hline
        Столбец & Значение & Тип & Ограничение \\ \hline
        \endhead

        \endfoot
        \endlastfoot

        id & Идентификатор & UUID & Не null, первичный ключ \\ \hline
        name & Название & Строка & Не null \\ \hline
        description & Описание & Строка & \\ \hline
        created\_at & Дата создания & Дата и время & Не null \\ \hline
        updated\_at & Дата обновления & Дата и время & Не null \\ \hline
        team\_id & Идентификатор \newline команды & UUID & Не null, внешний ключ \\ \hline
    \end{tabularx}
\end{center}

Для хранения топологии инфраструктуры в графовой базе данных спроектирована структура
узлов и связей. Данные сущности не имеют жесткой табличной схемы, но подчиняются правилам
модели графа свойств.

В графовой базе данных определен узел Component, атрибуты которого представлены в таблице
\ref{tbl:graph_node}.

\begin{center}
    \begin{tabularx}{\textwidth}{|p{3cm}|p{5.6cm}|p{1.8cm}|X|}
        \caption{\label{tbl:graph_node}Информация об атрибутах узла Component} \\ \hline
        Атрибут & Значение & Тип & Ограничение \\ \hline
        \endfirsthead

        \multicolumn{4}{l}{{Продолжение таблицы\ \thetable{}}} \\ \hline
        Атрибут & Значение & Тип & Ограничение \\ \hline
        \endhead

        \endfoot
        \endlastfoot

        id & Идентификатор & UUID & Не null, уникальный \\ \hline
        name & Название & Строка & Не null \\ \hline
        type & Тип компонента & Число & Не null \\ \hline
        description & Описание & Строка & \\ \hline
        system\_id & Идентификатор системы & UUID & Не null \\ \hline
    \end{tabularx}
\end{center}

Связь между узлами определяется направленным ребром DEPENDS\_ON. Информация об атрибутах
данного ребра представлена в таблице \ref{tbl:graph_edge}.

\begin{center}
    \begin{tabularx}{\textwidth}{|p{3cm}|p{5.6cm}|p{1.8cm}|X|}
        \caption{\label{tbl:graph_edge}Информация об атрибутах связи DEPENDS\_ON} \\ \hline
        Атрибут & Значение & Тип & Ограничение \\ \hline
        \endfirsthead

        \multicolumn{4}{l}{{Продолжение таблицы\ \thetable{}}} \\ \hline
        Атрибут & Значение & Тип & Ограничение \\ \hline
        \endhead

        \endfoot
        \endlastfoot

        id & Идентификатор & UUID & Не null, уникальный \\ \hline
        protocol & Тип протокола & Число & Не null \\ \hline
        severity & Уровень критичности & Число & Не null \\ \hline
    \end{tabularx}
\end{center}

На рисунке \ref{img:relational-dbml.pdf} представлена полная схема разрабатываемой реляционной базы данных.

\img{16cm}{relational-dbml.pdf}{Схема реляционной БД в нотации DBML}

\clearpage

\subsection{Ограничения целостности данных}

\subsubsection{Целостность таблиц и узлов}
Для обеспечения целостности реляционных таблиц каждая строка имеет уникальный
идентификатор, являющийся первичным ключом. Первичные ключи назначены полям id во всех
спроектированных таблицах. Для графовой базы данных на уровне СУБД устанавливается
ограничение уникальности на свойство id для всех узлов с меткой Component, что
предотвращает дублирование узлов графа.

\subsubsection{Целостность полей}
Для корректного функционирования системы накладываются ограничения на значения отдельных
полей. В таблице users поля username и email должны быть уникальными для
предотвращения создания дублирующихся учетных записей. В таблице invites поля code и 
target\_email также являются уникальным, а значение поля expiration\_date должно быть 
строго больше даты создания приглашения. В таблице it\_systems дата создания не может превышать текущее
системное время. В графовой базе данных атрибут type узла Component может принимать
только определенные значения: микросервис, база данных, брокер сообщений, внешний API.

\subsubsection{Целостность ссылок}
Обеспечение ссылочной целостности в реляционной базе данных реализуется с помощью внешних
ключей. Для каждого значения внешнего ключа в дочерней таблице должна существовать запись
с соответствующим первичным ключом в родительской таблице. В таблице \ref{tbl:fk}
представлены все внешние ключи реляционной схемы.

\begin{center}
    \begin{tabularx}{\textwidth}{|X|X|X|}
        \caption{\label{tbl:fk}Внешние ключи таблиц реляционной базы данных} \\ \hline
        Таблица & Внешний ключ & Ссылка на \newline родительскую таблицу \\ \hline
        \endfirsthead

        \multicolumn{3}{l}{{Продолжение таблицы\ \thetable{}}} \\ \hline
        Таблица & Внешний ключ & Ссылка на \newline родительскую таблицу \\ \hline
        \endhead

        \endfoot
        \endlastfoot

        users & teams\_id & teams (id) \\ \hline
        invites & teams\_id & teams (id) \\ \hline
        invites & activated\_by\_uid & users (id) \\ \hline
        it\_systems & teams\_id & teams (id) \\ \hline
    \end{tabularx}
\end{center}

\subsection{Ролевая модель на уровне приложения}
Для обеспечения безопасности и разграничения доступа к информации реализована ролевая
модель управления доступом на программном уровне. Базы данных не содержат индивидуальных
учетных записей для каждого пользователя системы. Вместо этого приложение подключается к
хранилищам данных с использованием единых учетных записей, а вся логика
проверки прав доступа и фильтрации запросов выполняется на уровне контроллеров и сервисов
разрабатываемого приложения.

В системе определены уровни доступа для пользователей, описанные в таблице~\ref{tbl:roles-desc}.

\begin{center}
    \begin{tabularx}{\textwidth}{|X|X|}
        \caption{\label{tbl:roles-desc}Описание ролей системы} \\ \hline
        Роль & Права \\ \hline
        \endfirsthead

        \multicolumn{2}{l}{{Продолжение таблицы\ \thetable{}}} \\ \hline
        Роль & Права \\ \hline
        \endhead

        \hline
        \endfoot
        \endlastfoot

        Незарегистрированный \newline пользователь & Имеет доступ исключительно к функциям
        регистрации, авторизации и валидации кода приглашения. \\ \hline
        
        Разработчик & Обладает правами на чтение данных о пользователях своей
        команды, а также на просмотр архитектурных схем систем, привязанных к его
        команде. Пользователь с данной ролью может инициировать запросы на визуализацию
        топологии, но не имеет прав на внесение изменений в структуру графа. \\ \hline

        SRE & Наследует все права разработчика, а также получает доступ к
        аналитическим функциям приложения. Данная роль позволяет запускать
        алгоритмические расчеты каскадных отказов, моделировать сценарии сбоев
        компонентов и запрашивать вычисление критических путей и единых точек отказа в графе. \\ \hline

        Архитектор & Обладает расширенными правами, включающими в себя возможности
        SRE-инженера, а также права на изменение топологии. Архитектор может отправлять
        запросы на создание, обновление и удаление систем, добавление новых узлов и
        настройку связей между ними в графовой базе данных. \\ \hline

        Администратор & Обладает полными правами в системе управления доступом.
        Пользователь с этой ролью может просматривать список всех зарегистрированных
        пользователей, создавать и удалять команды, генерировать коды приглашений. \\ \hline
    \end{tabularx}
\end{center}

\subsection[Проверка нахождения реляционной БД в 3НФ]{Проверка нахождения реляционной БД в 3НФ}

Для каждого отношения $R(U, F)$ разработанной реляционной базы данных проверка условий 3НФ 
выполняется следующим образом:

\subsubsection*{Классификация атрибутов отношения}
Все атрибуты из $U$ на основе их участия в функциональных зависимостях из множества $F$ делятся на 
три непересекающиеся группы:
\begin{itemize}
    \item $G_1$: атрибуты, встречающиеся только в левых частях ФЗ, либо вообще не 
    вошедшие в функциональные зависимости из множества $F$, эти атрибуты 
    входят в состав всех потенциальных ключей отношения;
    \item $G_2$: атрибуты, встречающиеся только в правых частях ФЗ, данные атрибуты 
    не входят в состав потенциальных ключей;
    \item $G_3$: промежуточные атрибуты, которые встречаются как в левых, так 
    и в правых частях ФЗ из множества $F$.
\end{itemize}

\subsubsection*{Определение потенциальных ключей}
Вычисляется замыкание множества атрибутов первой группы: $(G_1)^+$.
\begin{itemize}
    \item если $(G_1)^+ = U$, то множество $G_1$ является единственным минимальным потенциальным ключом отношения;
    \item если $(G_1)^+ \neq U$, то для получения потенциальных ключей базовая группа $G_1$ 
    последовательно комбинируется с подмножествами атрибутов из промежуточной группы $G_3$ до тех 
    пор, пока замыкание очередной комбинации не покроет все множество $U$.
\end{itemize}

\subsubsection*{Проверка требований 3НФ}
Отношение находится в 3НФ тогда и только тогда, когда для любой 
нетривиальной функциональной зависимости $X \to Y \in F^+$ выполняется хотя бы одно из следующих условий:
\begin{itemize}
    \item левая часть зависимости $X$ является суперключом отношения $R$;
    \item правая часть зависимости $Y$ состоит исключительно из первичных атрибутов, 
    входящих в состав хотя бы одного потенциального ключа.
\end{itemize}

\subsubsection{Отношение <<Команды>>}
Атрибуты:
\begin{equation*}
    U = \{\text{id}, \text{name}, \text{description}\}
\end{equation*}

Множество ФЗ:
\begin{equation*}
    F = \{ \text{id} \to \{\text{name}, \text{description}\} \}
\end{equation*}

Классификация атрибутов:
\begin{itemize}
    \item только в левой части: id; 
    \item только в правой части: {name, description};
    \item в обеих частях: $\emptyset$.
\end{itemize}

Нахождение потенциальных ключей:
\begin{equation*}
    (\text{id})^+ = \{\text{id}, \text{name}, \text{description}\} = U
\end{equation*}

Потенциальные ключи: id.

Проверка требований 3НФ: левые части всех функциональных зависимостей 
являются потенциальными ключами, следовательно, суперключами отношения <<Команды>>. Условие 3НФ выполняется.

\subsubsection{Отношение <<Пользователи>>}
Атрибуты:
\begin{equation*}
    U = \{\text{id}, \text{username}, \text{email}, \text{password\_hash}, \text{role}, \text{team\_id}\}
\end{equation*}

Множество ФЗ:
\begin{equation*}
    F = \{ \text{id} \to \{\text{username}, \text{email}, \text{password\_hash}, \text{role}, \text{team\_id}\},
\end{equation*}
\begin{equation*}
    \text{username} \to \{\text{id}, \text{email}, \text{password\_hash}, \text{role}, \text{team\_id}\},
\end{equation*}
\begin{equation*}
    \text{email} \to \{\text{id}, \text{username}, \text{password\_hash}, \text{role}, \text{team\_id}\} \}
\end{equation*}

Классификация атрибутов:
\begin{itemize}
    \item только в левой части: $\emptyset$;
    \item только в правой части: password\_hash, role, team\_id;
    \item в обеих частях: id, username, email.
\end{itemize}

Нахождение потенциальных ключей:
\begin{equation*}
    (\text{id})^+ = U
\end{equation*}
\begin{equation*}
    (\text{username})^+ = U
\end{equation*}
\begin{equation*}
    (\text{email})^+ = U
\end{equation*}

Потенциальные ключи: id, username и email.

Проверка требований 3НФ: левые части всех функциональных зависимостей являются потенциальными ключами,
следовательно, суперключами отношения <<Пользователи>>. Условие 3НФ выполняется.

\subsubsection{Отношение <<Приглашения>>}
Атрибуты:
\begin{align*}
    U &= \{\text{id}, \text{status}, \text{target\_email}, \text{role}, \text{code}, \\ 
    & \quad \text{expiration\_date}, \text{team\_id}, \text{activated\_by\_uid}\}    
\end{align*}

Множество ФЗ:
\begin{align*}
    F &= \{ \text{id} \to \{\text{status}, \text{target\_email}, \text{role}, \text{code}, \\ 
    & \quad \text{expiration\_date}, \text{team\_id}, \text{activated\_by\_uid}\}, \\
    & \text{code} \to \{\text{id}, \text{status}, \text{target\_email}, \text{role}, \\ 
    & \quad \text{expiration\_date}, \text{team\_id}, \text{activated\_by\_uid}\}, \\
    & \text{target\_email} \to \{\text{id}, \text{status}, \text{role}, \text{code}, \\ 
    & \quad \text{expiration\_date}, \text{team\_id}, \text{activated\_by\_uid}\} \}
\end{align*}

Классификация атрибутов:
\begin{itemize}
    \item только в левой части: $\emptyset$;
    \item только в правой части: status, role, expiration\_date, team\_id, \\ activated\_by\_uid;
    \item в обеих частях: id, code, target\_email.
\end{itemize}

Нахождение потенциальных ключей:
\begin{equation*}
    (\text{id})^+ = U
\end{equation*}
\begin{equation*}
    (\text{code})^+ = U
\end{equation*}
\begin{equation*}
    (\text{target\_email})^+ = U
\end{equation*}

Потенциальные ключи: id, code и target\_email.

Проверка требований 3НФ: левые части всех функциональных зависимостей являются потенциальными 
ключами отношения <<Приглашения>>. Условие 3НФ выполняется.

\subsubsection{Отношение <<Системы>>}
Атрибуты:
\begin{equation*}
    U = \{\text{id}, \text{name}, \text{description}, \text{created\_at}, \text{updated\_at}, \text{team\_id}\}
\end{equation*}

Множество ФЗ:
\begin{equation*}
    F = \{ \text{id} \to \{\text{name}, \text{description}, \text{created\_at}, \text{updated\_at}, \text{team\_id}\} \}
\end{equation*}

Классификация атрибутов:
\begin{itemize}
    \item только в левой части: id;
    \item только в правой части: name, description, created\_at, updated\_at, team\_id;
    \item В обеих частях: $\emptyset$.
\end{itemize}

Нахождение потенциальных ключей:
\begin{equation*}
    (\text{id})^+ = \{\text{id}, \text{name}, \text{description}, \text{created\_at}, \text{updated\_at}, \text{team\_id}\} = U
\end{equation*}

Потенциальный ключ: id.

Проверка требований 3НФ: левая часть единственной функциональной зависимости 
совпадает с потенциальным ключом отношения <<Системы>>. 
Транзитивные зависимости отсутствуют, условие 3НФ выполняется.

\subsubsection{Выводы по результатам анализа}
На основе проведенного анализа функциональных зависимостей и поиска потенциальных ключей 
для всех отношений реляционной базы данных сделаны следующие выводы:
\begin{itemize}
    \item для каждого из четырех отношений левая часть любой нетривиальной ФЗ является суперключом соответствующего отношения;
    \item в схеме базы данных отсутствуют транзитивные ФЗ непервичных атрибутов от потенциальных ключей.
\end{itemize}

Таким образом, спроектированная схема реляционной базы данных полностью удовлетворяет требованиям третьей нормальной формы.

\numberlesssubsection{Выводы по конструкторскому разделу}
В данном разделе была спроектирована структура хранения данных для системы анализа
отказоустойчивости. Описаны таблицы реляционной базы данных и структуры графовой
базы данных, определены связи между ними и ограничения целостности. Формализована ролевая
модель управления доступом на уровне приложения, обеспечивающая безопасность
взаимодействия с системой и разграничение функциональных возможностей пользователей.
